\section{Introduction}
An Agent receives the instruction \emph{process the next support ticket}. It reads the ticket, checks an order and approval, consults a model, and submits a refund. Encryption hides message contents from a network observer, but the contacted services, traffic sizes and timing remain visible. The tool's answer then becomes part of a model request, which creates a second observable representation of the same task. Suppose a defender adds cover traffic. A second question follows: does that cover end at a proxy, execute a read, or change business state? Figure~\ref{fig:chain} follows information and actions through this example.

\begin{figure*}[t]
\centering
\begin{tikzpicture}[x=1cm,y=1cm,font=\small,>=Latex]
\tikzset{box/.style={draw,rounded corners,align=center,minimum height=.85cm,text width=2.4cm},
 note/.style={align=center,font=\footnotesize,text width=3.1cm}}
\node[box,fill=blue!5] (u) at (1.3,1.15) {User\\``Process the next ticket''};
\node[box,fill=blue!5] (a) at (5.1,1.15) {Agent\\plan, check, act, retry};
\node[box,fill=green!6] (m) at (9.2,2.6) {Model service\\next action / streamed output};
\node[box,fill=green!6] (t) at (9.2,-.35) {Tool service\\ticket, order, refund};
\node[box,fill=orange!8] (d) at (14.1,1.15) {Specialist Agent\\optional delegation};
\draw[<->,thick] (u)--node[above,yshift=15pt]{client link C}(a);
\draw[<->,thick] (a.north east)--node[above,sloped]{model link M}(m.west);
\draw[<->,thick] (a.south east)--node[below,sloped]{tool link T}(t.west);
\draw[<->,thick] (a.east)--node[above]{coordination A}(d.west);
\node[note] at (2.0,-.35) {\textbf{Transit observer}\\lengths, gaps, connections;\\destinations when available};
\node[note] at (13.7,-.35) {\textbf{Recipient / active service}\\received arguments;\\optionally selected responses};
\node[draw=black!45,fill=black!3,align=left,text width=15.7cm,inner sep=6pt] at (7.9,-1.65) {
\textbf{Two protection questions:} Where does tool feedback appear next? What work does privacy cover execute before it is discarded? The original task and each observable service remain part of the protection decision.};
\end{tikzpicture}
\caption{Running support-ticket example and observation rights. A tool result returns to the Agent and can enter the next model request. Each encrypted link exposes its own traffic shape; a recipient additionally sees the request addressed to it. Concurrent tickets may share a connection on either link. The arrows represent application communication, not a one-to-one mapping from tool events to packets.}
\label{fig:chain}
\end{figure*}


\paragraph{Why systematize now}
Agent deployments combine remote tool protocols, including the Model Context Protocol (MCP), with streamed model APIs~\cite{mcp2025transports,rfc9113}. In under three years, studies have recovered research queries from contacted domains~\cite{jeong2026wra}, recognized tasks from model traffic~\cite{li2026ant,pouryousef2026traffic}, fingerprinted hosted Agents~\cite{agentprint2025} and attributed models from tool-call behavior~\cite{zhao2026agenttrace,ediga2026trace}. The first defenses are research decoys, relays and shaping~\cite{jeong2026wra,dangol2026workflow}. Most of this work is preprint-stage, and no shared evaluation standard exists. Website fingerprinting (WF) went through the same phase, and many early attack and defense results were later reassessed under realistic populations and adapted attackers~\cite{dyer2012peekaboo,juarez2014critical,mathews2023critical}. An SoK at this point can set evaluation expectations before results pile up that cannot be compared. Its corpus is necessarily small, so we make every conclusion's source dependence explicit.

\paragraph{The question}
Which traffic-analysis evaluation conditions carry over to Agent execution, and which break? We find that two assumptions behind WF-style defense evaluation fail: that one link is the unit of protection, and that cover traffic is free of side effects. Each failure leads to a concrete evaluation error. We first establish both failures from the literature, then measure their consequences in a controlled case. That case is the strongest test of the SoK's evaluation procedure. It applies the procedure to a cover defense that looks effective under the standard size-only evaluation, and the procedure shows that the defense fails.

\textbf{Finding 1: protection scope follows the task across service links.} WRA recovers research queries from tool destinations, while ANT and a cloud benchmark recognize task labels on model connections~\cite{jeong2026wra,li2026ant,pouryousef2026traffic}. A contacted tool additionally reads its arguments, and PrivacyTrace shows that these arguments raise attribute inference well above a schema prior~\cite{duan2026privacytrace}. Because tool results are forwarded into later model requests, a tool-side defense leaves the model link outside its scope. In our forwarding pilot, removing supplementary records at the first feedback step still leaves complete-execution model-view recognition at 44/120 for Qwen and 33/120 for Llama, where chance is 20/120 (Section~\ref{sec:live-study}). A defense claim must therefore name every view its attacker retains (Section~\ref{sec:authorization}).

\textbf{Finding 2: executed cover makes utility part of the privacy result.} Classical cover terminates at a tunnel endpoint~\cite{apthorpe2019stp}. An Agent decoy can execute queries or update state. The published research-decoy result lowers query recovery from .770 to .565, but its utility scores belong to a separate destination-visibility intervention~\cite{jeong2026wra}. Across the four mapped query/workflow privacy settings, adapted attacks appear in two and matched task utility in one, and no setting reports both (Section~\ref{sec:reporting}). These mapped records provide no setting with both adapted inference and matched task utility, so their joint evidence remains incomplete.

\textbf{Finding 3: executed cover carries the planner's timing signature.} An LLM Agent issues each honest tool call after a model decision, which takes hundreds of milliseconds. Cover inserted by a gateway, proxy or controller outside the model loop skips that wait. Padding hides sizes but cannot hide the missing decision time. We test this on 1,440 held-out MCP executions from two LLM planners (Section~\ref{sec:cover-study}). Under size-only observation, five cover policies cut task recognition from 113--118 to 61--82 of 120. An unsupervised gap rule then flags every decoy with no false removals. Stripping the flagged calls restores recognition to 107--116/120, and the attack still reaches 98--115/120 under 200\,ms of timestamp jitter. A generic classifier given the same timing features does not improve at all. The leak is found only by an attack designed around the execution mechanism. The fix follows from the cause. Drawing each decoy's dispatch gap and duration from the planner's own history returns the best attack to the size-only level, at 12--16\% added latency. Cover semantics decide the remaining trade-off: they move size-only recognition by at most ten tasks, while adverse writes destroy whole-state integrity in up to 113 of 120 fixtures.

\paragraph{Contributions}
\begin{itemize}
\item \textbf{An access-aware evidence map} of 44 settings from 15 Agent studies. It separates transit, recipient and active observers and binds each score to its target, population and supplied boundaries (Sections~\ref{sec:scope}--\ref{sec:taxonomy}, Tables~\ref{tab:deployment-access}--\ref{tab:attack-recipient}).
\item \textbf{A classical-to-Agent transfer matrix} that labels each established WF and traffic-analysis condition as retained, extended, untested or mismatched in access (Section~\ref{sec:correspondence}).
\item \textbf{Two execution-aware syntheses}, of feedback exposure and of cover termination, with a design space of six defense layers (Sections~\ref{sec:evidence}--\ref{sec:defenses}).
\item \textbf{A planner-latency attack on executed cover, with its countermeasure and cost}: a label-free removal attack that is robust to jitter, and timing-matched dispatch that neutralizes it, measured together with task, state and latency outcomes (Section~\ref{sec:cover-study}). A separate forwarding/concurrency pilot calibrates event-level evaluation (Section~\ref{sec:live-study}).
\item \textbf{A four-step evaluation procedure, an eight-field checklist and a source-withdrawal analysis} that shows which conclusions survive with only published sources (Sections~\ref{sec:reporting}--\ref{sec:agenda}).
\end{itemize}
